\section{评测体系演进：从 MBPP/HumanEval 到 SWE-Bench}

\subsection{早期代码评测的价值与失效}
MBPP、HumanEval 曾经推动了代码模型快速迭代，因为它们满足当时三项关键条件：难度合适、可自动判分、尚未广泛泄露。Sutton 直接指出，如今这类任务更像 ``code eval 的 MNIST''：仍可做基础对比，但难以代表真实软件工程能力。

\begin{knowledgebox}{Pass@K 的解释与部署含义}
Pass@K 衡量 ``采样 K 次至少有一次通过测试'' 的概率。它对评估多样性有价值，但并不总等价于产品体验：
\begin{itemize}
    \item IDE 单次补全更接近 Pass@1；
    \item Agent 带验证环节可以更接近 Pass@K；
    \item 因此指标是否有效，要看系统真实使用方式。
\end{itemize}
\end{knowledgebox}

\begin{warningbox}{指标错配会造成方向性误导}
如果产品场景只允许一次输出，却主要优化 Pass@K，团队可能会把系统训练成 ``样本池里有正确答案''，而不是 ``首答可用''。这种偏差在上线阶段会非常明显。
\end{warningbox}

\subsection{SWE-Bench 的突破}
SWE-Bench 把任务从 ``写短函数'' 推进到 ``在真实代码库里修真实 issue''：输入是 issue 描述与仓库上下文，输出是 patch，验证依赖已有测试套件。相比早期评测，它显著提升了真实度，也迫使研究者正视代码导航、上下文检索、多文件编辑等系统问题。

\begin{figure}[H]
\centering
\includegraphics[width=0.88\textwidth]{frame-02-swebench.jpg}
\caption{课程中对 SWE-Bench 任务形态与影响力的讲解画面 \protect\footnotemark}
\end{figure}
\footnotetext{视频画面时间区间：00:14:55--00:15:40。}

\begin{importantbox}{SWE-Bench 真正改变了什么}
它改变的不只是分数，而是研发范式。团队开始系统性投入：
\begin{itemize}
    \item 代码库级检索与导航；
    \item 执行反馈驱动修补；
    \item 多轮轨迹管理；
    \item Patch 生成与验证闭环。
\end{itemize}
\end{importantbox}

\subsection{SWE-Bench 的局限同样重要}
Sutton 讲得很坦率：SWE-Bench 的价值与局限并存。数据来源集中在少量项目和特定语言生态；issue 描述风格与 IDE 内真实人机交互不完全一致；历史提交存在 ``一个 commit 混合多个改动'' 的噪声；测试集并非完美 verifier，可能出现 ``过拟合测试但语义错误'' 的补丁。

\begin{table}[H]
\centering
\begin{tabular}{p{3.4cm}p{5.0cm}p{5.0cm}}
\toprule
\textbf{维度} & \textbf{SWE-Bench 贡献} & \textbf{潜在偏差} \\
\midrule
任务真实性 & 来自真实项目 issue/patch 链路 & 项目分布有限，领域覆盖不全 \\
验证方式 & 可自动运行测试，形成闭环评测 & 测试不能完整代表语义正确性 \\
语言输入 & 保留真实 issue 文风与上下文线索 & 与 IDE 内简短指令存在分布差异 \\
训练价值 & 可驱动 retrieval、editing、execution 设计 & 历史提交可能缠绕多种改动目标 \\
\bottomrule
\end{tabular}
\caption{SWE-Bench 的 ``强现实性'' 与 ``强偏差性'' 并存}
\end{table}

\subsection{评测寿命与持续再造}
讲座提出了一个非常关键的判断：\textbf{所有评测都有 shelf life}。起初难且未泄露，几年后常常变得易解且可被记忆。评测体系必须持续更新，否则系统会在过期目标上做高效优化。

\begin{importantbox}{一句可执行的方法论}
\textit{``If your numbers aren't good, you're hill climbing on the wrong thing.''}  
翻译成工程动作就是：先检查评测目标是否真实、是否仍有区分度，再做模型或 harness 调参。
\end{importantbox}

\subsection{本章小结}
从 MBPP/HumanEval 到 SWE-Bench，行业并不是 ``换了个榜单''，而是把代码智能从函数级推进到仓库级。与此同时，评测本身也必须被持续审计，否则高分不等于高可用。

